\documentclass[11pt]{article}
\usepackage{amsmath,amssymb,amsthm,geometry}
\geometry{margin=1in}\newtheorem{theorem}{Theorem}
\title{Eventual return of every fixed interior rank}
\author{AN04 closure campaign}\date{Version 1}
\begin{document}\maketitle
\paragraph{Leading system and normalization.} Fix $13/20\le\rho\le7/10$,
$\lambda=\rho^2$, $a=\rho K(\rho a)$, $b>0$ and $\zeta>a$.
Use the RMS join aligned law at $M=1,N=0$, with bottom half atom
at $a$ and continuous upper half the pushforward of
$(2/\pi)(1+u^2)^{-2}du$ on $u>0$ by $H^{-1}(bu)$.
Times are elapsed since this prescribed entry. IMP scope supplies
finite-interval characteristic evolution with locally finite first
absolute moments. The normalized rank parameter $f\in(0,1)$
is upper mass. Co-order is in this one field, and radial reaction
vanishes. $Y(0)>0$ implies $Y>0$ by $2Y'=-Y+\rho K_s$.

\begin{theorem}[Leading system]\label{an04ffpv1:return}
For every such evolution and every fixed rank $0<f<1$ there is a
finite time after which $x_f<\zeta$ permanently. Consequently both
the instantaneous inactive mass and the mass continuously inactive
from entry tend to zero. The parameter set with positive perpetual
fraction is empty, its complement is the entire stated domain,
and its relative boundary is empty. In fact $x_f(t)\to-\infty$ for every fixed interior rank.
\end{theorem}
\begin{proof}
Write $L=\lambda/2$. GOE ceiling (or RMS horizons) gives
\[
 (x_f)_+,(y_f)_+\le C_f e^{r_f t},\qquad
 r_f=\tfrac\lambda4(\sqrt{4+f^2}-f)<L.
\]
For the atom one can use any interior continuous rank above it by order.
This upper bound needs the positive initial mean and common-field order.

\emph{Moment at the moving receiver.} We restate the stronger
displacement information needed in addition to IMP's mass conclusion.
IMP budget gives $\max(\rho K_s,Y)<B e^{Lt}(1+t)^{-2/3}$.
At horizon $T$ take initial labels above
$R_T=1+\zeta+3B[(1+T)^{1/3}-1]$ and put
$w(t)=e^{Lt}(R_T-3B[(1+t)^{1/3}-1])$.
It satisfies $2w'=\lambda w-2B e^{Lt}(1+t)^{-2/3}$.
At $x=w,y\ge w$, $2x'\ge\rho K(\rho w)-\rho K_s$ since
$S\le K_s$. At $y=w,x\ge w$, $2y'\ge\rho K(\rho w)-Y$.
Both exceed $2w'$ since $\rho K(\rho w)>\lambda w$.
The two faces include the corner and the initial ranks meet both bounds.
Hence at $T$ their inputs exceed or equal $(1+\zeta)e^{LT}$.
Their mass is at least $c/(1+T)$ since the RMS join tail density
is asymptotic to $2b^3/(\pi x^4)$, is positive above $a$,
and $R_T\le C(1+T)^{1/3}$. This proof uses changing sets for
different horizons. It gives a tail moment at the fixed gate of at
least $D(t)=c(1+\zeta)e^{Lt}/(1+t)$ at every time.
Moreover, for $t$ sufficiently large depending on $f$,
$(1+\zeta)e^{Lt}>C_fe^{r_f t}$. Thus wherever $x_f\ge\zeta$,
the entire protected set lies above the \emph{current receiver},
including when that receiver is far beyond the gate.

\emph{Velocity until and after return.} If $\tilde x\ge x$ then
$F(\rho x,\rho\tilde x)=\phi(\rho x)+\rho\tilde x\Phi(\rho x)$.
Elsewhere $F\ge0$. For any $x\ge\zeta>0$ this gives
$S(x)\ge\rho\Phi(\rho x)\int_{\tilde x\ge x}\tilde x\,d\mu$.
(The positive gate matters when interpreting this as a moment bound.)
The doubled input velocity at the fixed rank thus obeys, whenever
$x_f\ge\zeta$ and $t$ is large,
\begin{equation}\label{an04ffpv1:drift}
 2x_f'\le\rho\phi(0)+\lambda\Phi(\rho x_f)
          [C_f e^{r_ft}-D(t)]\le-1 \quad\text{when }x_f\ge\zeta .
\end{equation}
Indeed $D(t)/e^{r_ft}\to\infty$ and $\Phi(\rho x_f)\ge1/2$.
The same formula with $x=\zeta$ verifies the proposed gate bound.
Starting at the finite time from which this inequality holds,
a rank above the gate must reach it in finite additional time.
The strict inward drift at equality places it below and forbids
recrossing; a rank already below cannot cross upward either.
Dominated convergence for the unit rank measure now proves the
instantaneous statement. The continuously inactive mass is bounded
above by that mass. The endpoint $f=0$ carries zero mass, and no
uniform return time in $f$ was used.

For the last assertion fix any finite lower threshold $z$, keeping
the protected set constructed with $\zeta>a$. For late $t$ it lies
above $x_f$ regardless of the receiver's sign. Restrict the nonnegative
overlap integral to this set alone, on which all donor inputs are
positive. This gives $S(x_f)\ge\rho\Phi(\rho x_f)D(t)$ even for
$x_f<0$. While $x_f\ge z$ use $\Phi(\rho x_f)\ge\Phi(\rho z)>0$
instead of $1/2$ in \eqref{an04ffpv1:drift}. The same crossing
argument yields permanent return below $z$. Apply this to each negative
integer. Thus the zero-fraction conclusion also extends to every finite
real gate if inactivity is defined by that cut.
\end{proof}

\paragraph{Leading-system positive route.}
The aligned wedge would require $\rho S(x)-Y\ge0$ on its diagonal;
at entry $\rho S\le\rho K_s<Y$ because $\rho K(\rho x)<x$ above $a$
and $x=y$ with positive upper mass. IMP mass still ensures survival
at least $c/(1+T)$ for every finite $T$. The theorem settles the
infinite-horizon question on both sides of the initial GOE threshold
and at equality. It gives no uniform decay rate and makes no weak-seed
or trained-flow assertion.
\end{document}
